% =====================================================================
% EXEMPLO PREENCHIDO (FICTÍCIO) — TCC NO FORMATO SBC
% Universidade Tiradentes (UNIT) — Curso de Ciência da Computação
% =====================================================================
% Este arquivo é uma versão de exemplo do modelo-tcc-sbc.tex, com todo
% o conteúdo preenchido com dados FICTÍCIOS (título, autor, orientador,
% resumo, seções e referências inventados) apenas para servir de
% referência visual de como o TCC preenchido deve ficar. NÃO é um
% trabalho real e não deve ser citado ou reaproveitado como fonte.
%
% A formatação (margens, fonte, coluna única, espaçamento) é idêntica
% à do modelo-tcc-sbc.tex em branco — veja aquele arquivo para os
% comentários explicando cada escolha de formatação.
% =====================================================================

\documentclass[a4paper,12pt]{article}

\usepackage[utf8]{inputenc}
\usepackage[T1]{fontenc}
\usepackage[brazilian]{babel}
\usepackage{times}
\usepackage[a4paper,top=3.5cm,left=3cm,right=3cm,bottom=2.5cm]{geometry}
\usepackage{graphicx}
\usepackage{url}
\usepackage{indentfirst}
\usepackage{titlesec}
\usepackage[bf,sf,footnotesize,indent]{caption}

\setlength{\parindent}{1.27cm}
\setlength{\parskip}{6pt}
\flushbottom

\titleformat{\section}{\fontsize{13}{15}\bfseries}{\thesection.}{1em}{}
\titleformat{\subsection}{\fontsize{12}{14}\bfseries}{\thesubsection.}{1em}{}
\titleformat{\subsubsection}{\fontsize{12}{14}\bfseries\itshape}{\thesubsubsection.}{1em}{}
\titlespacing*{\section}{0pt}{12pt}{6pt}
\titlespacing*{\subsection}{0pt}{10pt}{4pt}

\renewcommand{\captionfont}{\sffamily\footnotesize\bfseries}
\setlength{\captionmargin}{0.8cm}
\captionsetup{labelsep=period}

\newenvironment{resumo}{%
  \list{}{\advance\topsep by6pt\relax
    \leftmargin=0.8cm
    \labelwidth=0pt
    \listparindent=0pt
    \itemindent\listparindent
    \rightmargin\leftmargin}\item[\hskip\labelsep\bfseries\itshape Resumo.]\itshape}%
  {\endlist}

\renewenvironment{abstract}{%
  \list{}{\advance\topsep by6pt\relax
    \leftmargin=0.8cm
    \labelwidth=0pt
    \listparindent=0pt
    \itemindent\listparindent
    \rightmargin\leftmargin}\item[\hskip\labelsep\bfseries\itshape Abstract.]\itshape}%
  {\endlist}

\begin{document}
\pagestyle{empty}

\begin{center}
  {\fontsize{16}{20}\selectfont\bfseries Um Sistema de Recomendação de Disciplinas Eletivas Baseado em Aprendizagem de Máquina para Estudantes de Graduação}\\[4pt]
  {\fontsize{14}{18}\selectfont\bfseries A Machine Learning-Based Elective Course Recommendation System for Undergraduate Students}\\[12pt]

  {\normalsize\bfseries João Pedro Santos Almeida$^{1}$, Marina Costa Andrade$^{1}$, Ricardo Almeida Ferreira$^{2}$}\\[6pt]
  {\footnotesize\itshape Orientadora: Marina Costa Andrade (UNIT) -- Coorientador: Ricardo Almeida Ferreira (UFS)}\\[4pt]

  {\normalsize
    $^{1}$Curso de Ciência da Computação — Universidade Tiradentes (UNIT) -- Aracaju/SE\\{}
    $^{2}$Departamento de Computação — Universidade Federal de Sergipe (UFS) -- São Cristóvão/SE\\{}
    \footnotesize\texttt{\{joao.almeida,marina.andrade\}@souunit.com.br, ricardo.ferreira@ufs.br}
  }
\end{center}
\vspace{12pt}

\begin{resumo}
A escolha de disciplinas eletivas é uma etapa recorrente na vida acadêmica
que costuma ser feita sem apoio sistemático, levando estudantes a optarem
por disciplinas pouco alinhadas com seus interesses e objetivos de
carreira. Este trabalho propõe um sistema de recomendação de disciplinas
eletivas baseado em filtragem colaborativa, utilizando o histórico de
matrículas e as avaliações de desempenho de estudantes de cursos de
Ciência da Computação. Foram avaliados três algoritmos de aprendizagem de
máquina (KNN, SVD e um modelo baseado em fatoração de matrizes) sobre uma
base de dados simulada com 1.200 registros de matrícula. Os resultados
indicam que o modelo baseado em fatoração de matrizes obteve o melhor
desempenho, com precisão de 84,7\% nas recomendações do tipo top-5.
Conclui-se que sistemas de recomendação simples já trazem ganhos
relevantes de personalização para o planejamento acadêmico, servindo como
base para futuras integrações com sistemas de gestão acadêmica.
\end{resumo}
\noindent\textbf{Palavras-chave:} aprendizagem de máquina, sistema de recomendação, currículo acadêmico, filtragem colaborativa.

\vspace{6pt}

\begin{abstract}
Choosing elective courses is a recurring step in academic life that is
usually made without systematic support, leading students to pick courses
poorly aligned with their interests and career goals. This work proposes
an elective course recommendation system based on collaborative
filtering, using enrollment history and performance data from Computer
Science students. Three machine learning algorithms were evaluated (KNN,
SVD, and a matrix-factorization-based model) on a simulated dataset with
1,200 enrollment records. Results show that the matrix-factorization
model achieved the best performance, with 84.7\% top-5 recommendation
precision. We conclude that even simple recommendation systems bring
relevant personalization gains to academic planning, serving as a basis
for future integration with academic management systems.
\end{abstract}
\noindent\textbf{Keywords:} machine learning, recommender system, academic curriculum, collaborative filtering.

\vspace{18pt}

\section{Introdução}

A cada semestre, estudantes de graduação precisam escolher disciplinas
eletivas dentro de um conjunto amplo de opções, muitas vezes sem
informações suficientes sobre o conteúdo, a carga de trabalho ou a
aderência de cada disciplina aos seus interesses. Na prática, essa
escolha costuma ser guiada por recomendações informais de colegas ou pela
grade de horários disponível, e não por um critério que leve em conta o
histórico acadêmico do próprio estudante.

Este trabalho tem como objetivo geral propor e avaliar um sistema de
recomendação de disciplinas eletivas para estudantes de cursos de
Ciência da Computação, com base em técnicas de aprendizagem de máquina
aplicadas ao histórico acadêmico. Como objetivos específicos, busca-se
(i) construir uma base de dados simulada representativa do problema,
(ii) comparar diferentes algoritmos de recomendação e (iii) avaliar a
qualidade das recomendações geradas.

O restante deste trabalho está organizado da seguinte forma: a Seção 2
apresenta os trabalhos relacionados; a Seção 3 descreve a metodologia
adotada; a Seção 4 apresenta o desenvolvimento e os resultados obtidos; e
a Seção 5 traz as considerações finais.

\section{Trabalhos Relacionados}

Sistemas de recomendação aplicados ao contexto educacional têm sido
explorados sob diferentes abordagens. \cite{silva2022} propõem um sistema
baseado em regras para sugerir disciplinas com base em pré-requisitos
cursados, mas sem considerar o desempenho do estudante. \cite{oliveira2021}
utilizam filtragem colaborativa para recomendar materiais de estudo
complementares, obtendo ganhos de engajamento em ambientes virtuais de
aprendizagem. Já \cite{souza2023} comparam algoritmos de fatoração de
matrizes aplicados à recomendação de cursos de pós-graduação, reportando
resultados superiores aos de abordagens baseadas em conteúdo.

Diferentemente desses trabalhos, esta pesquisa foca especificamente em
disciplinas eletivas de graduação em Computação e compara três técnicas
distintas sob as mesmas condições experimentais, o que permite uma
análise mais direta de qual abordagem é mais adequada a esse cenário.

\section{Metodologia}

Foi construída uma base de dados simulada com 1.200 registros de
matrícula, representando 150 estudantes fictícios e 24 disciplinas
eletivas do curso de Ciência da Computação. Cada registro associa um
estudante, uma disciplina cursada e uma nota de avaliação de satisfação
(escala de 1 a 5), usada como proxy de interesse.

Três algoritmos foram implementados em Python, utilizando a biblioteca
scikit-surprise: (i) K-Nearest Neighbors (KNN) com similaridade de
cosseno; (ii) Singular Value Decomposition (SVD); e (iii) um modelo de
fatoração de matrizes com regularização L2. A base foi dividida em 80\%
para treino e 20\% para teste, e a avaliação utilizou a métrica de
precisão top-5, que mede a proporção de disciplinas relevantes entre as
cinco primeiras recomendações geradas para cada estudante.

\section{Desenvolvimento e Resultados}

Os três algoritmos foram treinados sobre a base descrita na Seção 3 e
avaliados no conjunto de teste. A Tabela \ref{tab:resultados} apresenta a
precisão top-5 obtida por cada algoritmo.

\begin{table}[ht]
  \centering
  \caption{Precisão top-5 por algoritmo avaliado.}
  \label{tab:resultados}
  \begin{tabular}{lc}
    \hline
    Algoritmo & Precisão top-5 \\ \hline
    KNN (similaridade de cosseno) & 71,3\% \\
    SVD & 79,5\% \\
    Fatoração de matrizes (L2) & 84,7\% \\ \hline
  \end{tabular}
\end{table}

O modelo de fatoração de matrizes apresentou o melhor desempenho entre os
três, superando o SVD em 5,2 pontos percentuais. Uma hipótese para essa
diferença é a capacidade do modelo de capturar padrões latentes de
afinidade entre estudantes com trajetórias acadêmicas parecidas, mesmo
quando o número de disciplinas em comum entre eles é pequeno --
característica frequente em bases esparsas como a utilizada neste
trabalho.

A Figura \ref{fig:arquitetura} ilustra a arquitetura geral do protótipo
desenvolvido, dividida em três módulos: coleta e pré-processamento dos
dados de matrícula, treinamento do modelo de recomendação e geração das
recomendações apresentadas ao estudante.

\begin{figure}[ht]
  \centering
  \fbox{\parbox{0.85\linewidth}{\centering\vspace{1.2cm}
  \small (espaço reservado para a figura da arquitetura do protótipo)
  \vspace{1.2cm}}}
  \caption{Arquitetura geral do protótipo de recomendação.}
  \label{fig:arquitetura}
\end{figure}

\section{Conclusão}

Este trabalho propôs e avaliou um sistema de recomendação de disciplinas
eletivas para estudantes de Ciência da Computação, comparando três
algoritmos de aprendizagem de máquina sobre uma base de dados simulada.
O modelo de fatoração de matrizes obteve o melhor resultado, com 84,7\%
de precisão top-5, indicando que técnicas relativamente simples de
filtragem colaborativa já são capazes de gerar recomendações relevantes
para o planejamento acadêmico.

Como limitação, destaca-se o uso de dados simulados em vez de dados reais
de matrícula, o que pode não capturar toda a variabilidade do
comportamento real dos estudantes. Como trabalhos futuros, sugere-se (i)
validar o sistema com dados reais e anonimizados de uma instituição de
ensino, e (ii) integrar o protótipo a um sistema acadêmico em produção
para avaliação com usuários reais.

\section*{Declaração de Uso de Inteligência Artificial}

Em conformidade com a Portaria CNPq nº 2.664/2026 e o Art. 12º do Manual
de Normas de TCC da Universidade Tiradentes, os autores declaram que a
ferramenta de IA generativa ChatGPT (OpenAI, modelo GPT-4) foi utilizada
exclusivamente nas etapas de \emph{revisão ortográfica e gramatical} do
texto em português e de \emph{tradução} do Resumo e das palavras-chave
para o inglês (Abstract e Keywords). Nenhuma ferramenta de IA generativa
foi utilizada na geração de código-fonte, na análise dos dados
experimentais, na obtenção dos resultados reportados na Tabela
\ref{tab:resultados} ou na redação do conteúdo técnico e científico do
trabalho, cuja autoria e responsabilidade integral permanecem dos
autores.

\begin{thebibliography}{9}

\bibitem{silva2022}
Silva, C. R.
\newblock Recomendação de disciplinas baseada em regras de pré-requisito.
\newblock \emph{Anais do Workshop de Informática na Educação}, p. 45--54, 2022.

\bibitem{oliveira2021}
Oliveira, F. A.; Souza, L. M.
\newblock Filtragem colaborativa para recomendação de materiais de estudo
em ambientes virtuais de aprendizagem.
\newblock \emph{Revista Brasileira de Informática na Educação}, v. 29,
n. 2, p. 112--130, 2021.

\bibitem{souza2023}
Souza, T. P.; Lima, R. B.; Costa, H. N.
\newblock Fatoração de matrizes aplicada à recomendação de cursos de
pós-graduação: um estudo comparativo.
\newblock \emph{Anais do Simpósio Brasileiro de Sistemas de Informação},
p. 201--210, 2023.

\end{thebibliography}

\end{document}
